% ==============================================================================
% Paper 84 (EN): The Fontaine-Mazur Conjecture:
%                p-adic Galois Representations and Modularity
% Author: Michael Fuhrmann
% Project: Specialist Mathematics Project, Build 167
% Date: 2026-03-12
% ==============================================================================

\documentclass[12pt,a4paper]{amsart}

\usepackage{amsmath,amssymb,amsthm,mathtools,enumitem,booktabs,array,hyperref}
\usepackage[utf8]{inputenc}
\usepackage[T1]{fontenc}
\usepackage{geometry}
\geometry{margin=2.5cm}

% --------------------------------------------------------------------------
% Theorem environments
% --------------------------------------------------------------------------
\newtheorem{theorem}{Theorem}[section]
\newtheorem{lemma}[theorem]{Lemma}
\newtheorem{corollary}[theorem]{Corollary}
\newtheorem{proposition}[theorem]{Proposition}
\theoremstyle{definition}
\newtheorem{definition}[theorem]{Definition}
\newtheorem{example}[theorem]{Example}
\theoremstyle{remark}
\newtheorem{remark}[theorem]{Remark}
\newtheorem{conjecture}[theorem]{Conjecture}

% --------------------------------------------------------------------------
% Macros
% --------------------------------------------------------------------------
\newcommand{\Gal}{\mathrm{Gal}}
\newcommand{\GL}{\mathrm{GL}}
\newcommand{\SL}{\mathrm{SL}}
\newcommand{\Qp}{\mathbb{Q}_p}
\newcommand{\Zp}{\mathbb{Z}_p}
\newcommand{\Qbar}{\overline{\mathbb{Q}}}
\newcommand{\GQ}{G_{\mathbb{Q}}}
\newcommand{\BdR}{B_{\mathrm{dR}}}
\newcommand{\Bcris}{B_{\mathrm{cris}}}
\newcommand{\Bst}{B_{\mathrm{st}}}
\newcommand{\DdR}{D_{\mathrm{dR}}}
\newcommand{\Dcris}{D_{\mathrm{cris}}}
\newcommand{\Dst}{D_{\mathrm{st}}}
\newcommand{\HT}{\mathrm{HT}}
\newcommand{\WD}{\mathrm{WD}}
\newcommand{\Sel}{\mathrm{Sel}}
\newcommand{\Frob}{\mathrm{Frob}}
\newcommand{\rec}{\mathrm{rec}}
\newcommand{\Z}{\mathbb{Z}}
\newcommand{\Q}{\mathbb{Q}}
\newcommand{\R}{\mathbb{R}}
\newcommand{\C}{\mathbb{C}}
\newcommand{\F}{\mathbb{F}}
\newcommand{\N}{\mathbb{N}}
\newcommand{\HH}{\mathcal{H}}

\hypersetup{
  colorlinks=true,
  linkcolor=blue,
  citecolor=blue,
  urlcolor=blue,
  pdftitle={The Fontaine-Mazur Conjecture: p-adic Galois Representations and Modularity},
  pdfauthor={Michael Fuhrmann}
}

% ==============================================================================
\begin{document}
% ==============================================================================

\title{The Fontaine-Mazur Conjecture:\\
  $p$-adic Galois Representations and Modularity}

\author{Michael Fuhrmann}
\address{Specialist Mathematics Project, Build 167}
\email{specialist-maths@project.local}
\date{2026-03-12}

\begin{abstract}
The Fontaine-Mazur conjecture, formulated in 1995, predicts that every irreducible
geometric $p$-adic Galois representation
$\rho\colon \GQ \to \GL_n(\Qp)$
arises from the cohomology of a smooth proper variety over $\Q$, and in the two-dimensional
case is modular in the sense that it corresponds to a classical automorphic form.
We survey the foundational theory of $p$-adic Hodge theory developed by Fontaine,
including the period rings $\BdR$, $\Bcris$, $\Bst$ and the associated admissible
representations (de Rham, crystalline, semistable, Hodge-Tate).
We then state the Fontaine-Mazur conjecture precisely, describe its connection to the
Langlands program, and review the major proven cases:
Wiles' theorem for $n=2$ weight $2$ (with Breuil-Conrad-Diamond-Taylor completing the
proof for all elliptic curves), Kisin's (2009) result on potentially semistable
representations of Hodge-Tate weights $\geq 1$, Emerton's completed cohomology approach,
and Taylor's theorem on potential modularity.
For $n \geq 3$ the conjecture remains wide open.
\end{abstract}

\maketitle

\tableofcontents

% ==============================================================================
\section{Introduction}
% ==============================================================================

One of the central problems in modern number theory is to understand all
$\ell$-adic Galois representations
\[
  \rho \colon \GQ \;:=\; \Gal(\Qbar/\Q) \;\longrightarrow\; \GL_n(\overline{\Q}_\ell),
\]
particularly to classify which of them \emph{come from geometry} (i.e.\ appear as
sub-quotients of the étale cohomology of smooth proper varieties over $\Q$) and which
are \emph{automorphic} (i.e.\ correspond to automorphic representations of $\GL_n$ via
the Langlands correspondence).

The Fontaine-Mazur conjecture \cite{FM1995} is a precise prediction of this kind for
$p$-adic Galois representations: an irreducible, geometric, odd (when $n=2$)
$p$-adic Galois representation should be modular / automorphic.

\subsection{Historical context}
The study of $p$-adic Galois representations was transformed in the 1980s and 1990s by
Fontaine's systematic development of $p$-adic Hodge theory
\cite{Fontaine1979,Fontaine1982,Fontaine1994}.
Fontaine introduced a hierarchy of ``nice'' ($p$-adic Hodge-theoretic) representations:
Hodge-Tate $\supset$ de Rham $\supset$ semistable $\supset$ crystalline,
and associated to each a module with additional structure (filtered $\phi$-module,
$(\phi,N)$-module, etc.).

In 1995, Fontaine and Mazur \cite{FM1995} formulated a sweeping conjecture asserting that
every irreducible \emph{geometric} representation (i.e.\ de Rham at $p$ and unramified
almost everywhere) arises from geometry and is hence automorphic via the Langlands program.

The $n=2$ case was established in its most important instances by Wiles \cite{Wiles1995}
and Taylor-Wiles \cite{TW1995} (semi-stable elliptic curves, later extended to all
elliptic curves by Breuil-Conrad-Diamond-Taylor \cite{BCDT2001}),
and in much greater generality by Kisin \cite{Kisin2009} and Emerton \cite{Emerton2010}.

\subsection{Structure of this paper}
Section~\ref{sec:padichg} reviews $p$-adic Hodge theory and Fontaine's period rings.
Section~\ref{sec:admissible} describes admissible representations.
Section~\ref{sec:galrep} sets up the general theory of $p$-adic Galois representations.
Section~\ref{sec:fm-statement} states the Fontaine-Mazur conjecture.
Section~\ref{sec:langlands} explains the connection to the Langlands program.
Section~\ref{sec:wiles} reviews Wiles' theorem and modularity of elliptic curves.
Section~\ref{sec:kisin} describes Kisin's work on potentially semistable representations.
Section~\ref{sec:emerton} covers Emerton's completed cohomology.
Section~\ref{sec:taylor} discusses Taylor's potential modularity.
Section~\ref{sec:openproblems} lists open problems.
Section~\ref{sec:conclusion} concludes.

% ==============================================================================
\section{$p$-adic Hodge Theory: Period Rings}\label{sec:padichg}
% ==============================================================================

Let $p$ be a prime, and let $K/\Qp$ be a finite extension with absolute Galois group
$G_K = \Gal(\overline{K}/K)$. Fontaine constructed a hierarchy of period rings that allow
one to attach linear-algebraic data to $p$-adic representations.

\subsection{The de Rham period ring $\BdR$}

\begin{definition}[Fontaine \cite{Fontaine1982}]
  Let $\widetilde{\mathbf{E}}^+ = \varprojlim_{x \mapsto x^p} \mathcal{O}_{\overline{K}}/p$
  be the tilt of $\mathcal{O}_{\overline{K}}$. The ring $A_{\inf}(\mathcal{O}_{\overline{K}})$
  is the Witt vector ring $W(\widetilde{\mathbf{E}}^+)$. There is a surjective ring
  homomorphism $\theta\colon A_{\inf} \to \mathcal{O}_{\mathbb{C}_p}$ whose kernel is a
  principal ideal generated by a non-zero-divisor $\xi$.

  The \emph{de Rham period ring} is
  \[
    \BdR \;:=\; \varprojlim_{n} A_{\inf}[\tfrac{1}{p}] / (\ker\theta)^n,
  \]
  a complete discrete valuation field with residue field $\mathbb{C}_p$ and uniformizer
  $t = \log([\varepsilon])$, where $\varepsilon = (1, \zeta_p, \zeta_{p^2},\ldots)$ is a
  compatible system of $p$-power roots of unity.
\end{definition}

The ring $\BdR$ has a filtration $\mathrm{Fil}^i \BdR = t^i \BdR^+$ and a $G_K$-action;
the invariants satisfy $\BdR^{G_K} = K$.

\subsection{The crystalline period ring $\Bcris$}

\begin{definition}
  The \emph{crystalline period ring} $\Bcris$ is the $p$-adic completion of
  \[
    A_{\mathrm{cris}} = \left\{ \sum_{n=0}^\infty a_n \frac{\xi^n}{n!} \,:\, a_n \in A_{\inf}, \; a_n \to 0 \right\},
  \]
  inverted at $p$. It carries a Frobenius $\phi$ (induced from Witt vectors), a filtration
  (induced from $\BdR$), and a $G_K$-action. Its invariants satisfy $\Bcris^{G_K} = K_0$
  (the maximal unramified subfield of $K$).
\end{definition}

\subsection{The semistable period ring $\Bst$}

\begin{definition}
  Fix a uniformizer $\pi \in K$ and let $\ell = \log([\widetilde{\pi}]) \in \Bcris$,
  where $\widetilde{\pi} = (\pi, \pi^{1/p},\ldots)$. The \emph{semistable period ring} is
  \[
    \Bst \;:=\; \Bcris[\ell].
  \]
  It carries the Frobenius $\phi$ (extending from $\Bcris$), a monodromy operator
  $N\colon \Bst \to \Bst$ defined by $N(\ell) = -1$, $N|_{\Bcris} = 0$, and a
  $G_K$-action. We have $\Bcris \subset \Bst \subset \BdR$.
\end{definition}

\begin{remark}
  The inclusion $\Bcris \subsetneq \Bst$ is strict whenever $K/\Qp$ is ramified.
  For the cyclotomic character $\chi_{\mathrm{cyc}}\colon G_K \to \Zp^\times$, the
  element $t \in \BdR^+$ satisfies $g \cdot t = \chi_{\mathrm{cyc}}(g) \cdot t$.
\end{remark}

% ==============================================================================
\section{Admissible Representations}\label{sec:admissible}
% ==============================================================================

Let $V$ be a finite-dimensional $\Qp$-vector space with a continuous $G_K$-action
(a $p$-adic representation of $G_K$).

\begin{definition}
  \begin{enumerate}[label=(\roman*)]
    \item $V$ is \emph{Hodge-Tate} if $V \otimes_{\Qp} \mathbb{C}_p \cong \bigoplus_i \mathbb{C}_p(i)^{h_i}$
      as $G_K$-modules, where $\mathbb{C}_p(i)$ denotes the $i$-th Tate twist and $h_i \geq 0$.
    \item $V$ is \emph{de Rham} if $\DdR(V) := (V \otimes_{\Qp} \BdR)^{G_K}$ has
      $K$-dimension equal to $\dim_{\Qp} V$.
    \item $V$ is \emph{semistable} if $\Dst(V) := (V \otimes_{\Qp} \Bst)^{G_K}$
      has $K_0$-dimension equal to $\dim_{\Qp} V$.
    \item $V$ is \emph{crystalline} if $\Dcris(V) := (V \otimes_{\Qp} \Bcris)^{G_K}$
      has $K_0$-dimension equal to $\dim_{\Qp} V$.
  \end{enumerate}
\end{definition}

\begin{theorem}[Fontaine {\cite[Theorem A]{Fontaine1994}}]
  There is a strict inclusion of full subcategories of $p$-adic representations of $G_K$:
  \[
    \{ \text{crystalline} \} \;\subsetneq\; \{ \text{semistable} \} \;\subsetneq\;
    \{ \text{de Rham} \} \;\subsetneq\; \{ \text{Hodge-Tate} \}.
  \]
  Moreover, each inclusion is strict.
\end{theorem}

\begin{theorem}[Berger \cite{Berger2002}, $C_{\mathrm{dR}}$-conjecture = Fontaine's conjecture]
  Every de Rham representation of $G_K$ is potentially semistable, i.e.\ becomes
  semistable over some finite extension of $K$.
\end{theorem}

This theorem, proved by Berger \cite{Berger2002} using $p$-adic differential equations
(work of Colmez-Fontaine, Fontaine-Wach, and Berger), is fundamental for the Fontaine-Mazur
conjecture.

\begin{definition}
  The \emph{Hodge-Tate weights} of a de Rham representation $V$ of dimension $n$ are
  the integers $h_1 \leq h_2 \leq \cdots \leq h_n$ such that
  $\mathrm{gr}^i \DdR(V) \neq 0$ for each $h_j$.
\end{definition}

\begin{example}
  The cyclotomic character $\chi_{\mathrm{cyc}}\colon G_{\Qp} \to \Zp^\times$ is
  crystalline of Hodge-Tate weight $1$. The $p$-adic Tate module $T_p E$ of an elliptic
  curve $E/\Q_p$ with good reduction is crystalline of Hodge-Tate weights $\{0, 1\}$.
\end{example}

% ==============================================================================
\section{$p$-adic Galois Representations: Global Theory}\label{sec:galrep}
% ==============================================================================

We now work globally with $\GQ = \Gal(\Qbar/\Q)$.

\begin{definition}
  A continuous representation
  $\rho\colon \GQ \to \GL_n(\Qp)$
  is called:
  \begin{enumerate}[label=(\roman*)]
    \item \emph{unramified at $\ell \neq p$} if $\rho|_{I_\ell}$ is trivial, where
      $I_\ell$ is the inertia group at $\ell$;
    \item \emph{geometric} (in the sense of Fontaine-Mazur) if it is de Rham at $p$
      and unramified at almost all primes $\ell$;
    \item \emph{of geometric origin} if it appears as a sub-quotient of
      $H^i_{\text{\'et}}(X_{\Qbar}, \Qp)(j)$ for some smooth proper $X/\Q$, integer $i$,
      and Tate twist $j$.
  \end{enumerate}
\end{definition}

\begin{example}[Tate module]
  Let $E/\Q$ be an elliptic curve. The $p$-adic Tate module
  $V_p E = T_p E \otimes_{\Zp} \Qp$
  gives an irreducible $2$-dimensional $p$-adic Galois representation. It is:
  \begin{itemize}
    \item de Rham at $p$ (crystalline if $E$ has good reduction at $p$, semistable if
      multiplicative reduction);
    \item unramified at all $\ell \nmid p \cdot N_E$ (where $N_E$ is the conductor);
    \item hence geometric.
  \end{itemize}
\end{example}

\begin{example}[Galois representation attached to a modular form]
  Let $f = \sum_{n=1}^\infty a_n q^n \in S_k(\Gamma_0(N), \chi)$ be a normalized newform
  of weight $k \geq 2$. Deligne \cite{Deligne1971} constructed a $2$-dimensional
  $\lambda$-adic representation
  \[
    \rho_f \colon \GQ \to \GL_2(\mathcal{O}_{f,\lambda})
  \]
  (where $\mathcal{O}_{f,\lambda}$ is a $p$-adic ring of integers) that is:
  \begin{itemize}
    \item unramified at all $\ell \nmid Np$;
    \item de Rham at $p$ with Hodge-Tate weights $\{0, k-1\}$;
    \item irreducible (for $f$ not CM, by Ribet).
  \end{itemize}
\end{example}

% ==============================================================================
\section{The Fontaine-Mazur Conjecture}\label{sec:fm-statement}
% ==============================================================================

We can now state the main conjecture.

\begin{conjecture}[Fontaine-Mazur \cite{FM1995}]\label{conj:FM}
  Let $\rho\colon \GQ \to \GL_n(\Qp)$ be an irreducible, geometric $p$-adic Galois
  representation. Then:
  \begin{enumerate}[label=(\alph*)]
    \item $\rho$ is of geometric origin: it appears in the $p$-adic étale cohomology of a
      smooth proper variety over $\Q$.
    \item \textbf{(Modularity)} For $n=2$, if additionally $\rho$ is odd (i.e.\
      $\det\rho(\mathrm{complex\;conjugation}) = -1$), then $\rho \cong \rho_f$ for some
      classical modular form $f$.
    \item More generally, $\rho$ corresponds to an automorphic representation of $\GL_n$
      via the global Langlands correspondence.
  \end{enumerate}
\end{conjecture}

\begin{remark}
  Part (a) is sometimes called the \emph{motivic} aspect and part (b)/(c) the
  \emph{automorphic} aspect of the conjecture. The two are not obviously equivalent: a
  priori, geometric origin is a stronger condition than being geometric (de Rham + almost
  everywhere unramified), though in practice they coincide in all known cases.
\end{remark}

\begin{remark}
  The ``odd'' condition for $n=2$ is necessary: the cyclotomic character
  $\chi_{\mathrm{cyc}}$ itself is geometric but even, and corresponds to a Hecke
  Grössencharacter, not a classical holomorphic modular form.
\end{remark}

\begin{conjecture}[Fontaine-Mazur, $n=2$ case]\label{conj:FM2}
  Every odd, irreducible, geometric $2$-dimensional $p$-adic Galois representation
  $\rho\colon \GQ \to \GL_2(\Qp)$
  is modular: there exists a normalized cuspidal newform $f$ of weight $k \geq 1$ such
  that $\rho \cong \rho_f \otimes \Qp$.
\end{conjecture}

% ==============================================================================
\section{Connection to the Langlands Program}\label{sec:langlands}
% ==============================================================================

The Fontaine-Mazur conjecture is deeply embedded in the broader Langlands program, which
predicts a correspondence between:
\begin{itemize}
  \item \emph{Galois side}: $n$-dimensional representations of $\GQ$ (or more generally of
    the Weil-Deligne group $W_{\Q_\ell}$);
  \item \emph{Automorphic side}: automorphic representations $\pi$ of $\GL_n(\mathbb{A}_\Q)$,
    where $\mathbb{A}_\Q$ is the adele ring.
\end{itemize}

\begin{definition}[Local Langlands correspondence at $\ell \neq p$]
  For each prime $\ell$, there is a local Langlands correspondence (proved by Harris-Taylor
  \cite{HT2001} and Henniart \cite{Henniart2000} for $\GL_n$):
  \[
    \rec \colon \{ \text{irred.\ smooth rep.\ of } \GL_n(\Q_\ell) \} \;\longleftrightarrow\;
    \{ \text{irred.\ deg.\ $n$ WD rep.\ of } W_{\Q_\ell} \}.
  \]
\end{definition}

\begin{definition}[Weil-Deligne representation]
  A \emph{Weil-Deligne representation} of $W_K$ (the Weil group of $K$) is a pair
  $(r, N)$ where $r\colon W_K \to \GL_n(\C)$ is a representation and
  $N\colon \C^n \to \C^n$ is a nilpotent endomorphism satisfying
  $r(\sigma) N r(\sigma)^{-1} = |\sigma| N$.
\end{definition}

The local Langlands at $p$ (i.e.\ the $p$-adic local Langlands) is far more subtle and
is currently only fully understood for $\GL_2(\Qp)$:

\begin{theorem}[$p$-adic local Langlands for $\GL_2(\Qp)$, Colmez \cite{Colmez2010}]
  There is a $p$-adic local Langlands correspondence for $\GL_2(\Qp)$: a bijection
  (functorial, compatible with twists and duality) between
  \begin{itemize}
    \item absolutely irreducible $2$-dimensional $p$-adic representations of $G_{\Qp}$, and
    \item absolutely irreducible unitary admissible Banach space representations of
      $\GL_2(\Qp)$ with a central character.
  \end{itemize}
\end{theorem}

The Fontaine-Mazur conjecture for $n=2$ thus fits into the framework:
a global geometric $\rho$ should correspond (via the global Langlands correspondence)
to an automorphic representation $\pi$ of $\GL_2(\mathbb{A}_\Q)$ whose local component
at $p$ is given by Colmez's functor $\check{V} \mapsto \Pi(\check{V})$.

% ==============================================================================
\section{Wiles' Theorem and Modularity of Elliptic Curves}\label{sec:wiles}
% ==============================================================================

The first major breakthrough was Wiles' proof of the Shimura-Taniyama conjecture for
semistable elliptic curves \cite{Wiles1995,TW1995}.

\begin{theorem}[Wiles 1995, completed by Breuil-Conrad-Diamond-Taylor \cite{BCDT2001}]
  Every elliptic curve $E/\Q$ is modular: there exists a normalized cuspidal newform
  $f \in S_2(\Gamma_0(N))$ such that
  $L(E,s) = L(f,s)$,
  or equivalently, $V_p E \cong \rho_f$ as $\GQ$-representations.
\end{theorem}

\begin{remark}
  In terms of the Fontaine-Mazur conjecture, this proves Conjecture~\ref{conj:FM}(b) for
  all $\rho = V_p E$ coming from elliptic curves: these are $2$-dimensional, odd,
  irreducible (for $E$ without complex multiplication, by Serre), geometric with
  Hodge-Tate weights $\{0,1\}$.
\end{remark}

The proof proceeds through:
\begin{enumerate}[label=(\roman*)]
  \item \textbf{Residual modularity}: For $p=3$ or $p=5$, the residual representation
    $\bar{\rho}_E \colon \GQ \to \GL_2(\F_p)$ is either reducible or known to be modular
    (using Langlands-Tunnell for $\GL_2(\F_3)$, or a base-change trick for $p=5$).
  \item \textbf{Modularity lifting (R=T)}: A deformation-theoretic argument shows that
    the universal deformation ring $R$ of $\bar{\rho}_E$ maps isomorphically onto the
    Hecke algebra $\mathbb{T}$, hence any deformation lifting $\bar{\rho}_E$ that is
    geometric comes from a modular form.
\end{enumerate}

\begin{theorem}[Taylor-Wiles method \cite{TW1995}]
  Let $p$ be an odd prime and $\bar{\rho}\colon \GQ \to \GL_2(\F_p)$ be an absolutely
  irreducible, odd representation that is modular. Suppose $\rho\colon \GQ \to \GL_2(\Zp)$
  is a geometric deformation of $\bar{\rho}$ with prescribed Hodge-Tate weights.
  Then $\rho$ is modular.
\end{theorem}

% ==============================================================================
\section{Kisin's Work: Potentially Semistable Representations}\label{sec:kisin}
% ==============================================================================

Kisin \cite{Kisin2009} made a decisive advance towards the full $n=2$ Fontaine-Mazur
conjecture.

\begin{definition}[Potentially semistable]
  A $p$-adic representation $V$ of $G_K$ is \emph{potentially semistable} if there
  exists a finite extension $K'/K$ such that $V|_{G_{K'}}$ is semistable.
  By Berger's theorem, this is equivalent to $V$ being de Rham.
\end{definition}

\begin{theorem}[Kisin \cite{Kisin2009}]\label{thm:kisin}
  Let $p > 2$ and let $\bar{\rho}\colon \GQ \to \GL_2(\F_p)$ be an odd, absolutely
  irreducible representation. Suppose:
  \begin{enumerate}[label=(\roman*)]
    \item $\bar{\rho}|_{G_{\Qp}}$ is not isomorphic to
      $\begin{pmatrix} \chi_1 & * \\ 0 & \chi_2 \end{pmatrix}$
      with $\chi_1\chi_2^{-1} = \bar{\chi}_{\mathrm{cyc}}$;
    \item $\bar{\rho}|_{G_{\Q(\zeta_p)}}$ is absolutely irreducible;
    \item $\bar{\rho}$ is modular.
  \end{enumerate}
  Then any odd, irreducible $\rho\colon \GQ \to \GL_2(\Qp)$ with $\bar{\rho}_\rho = \bar{\rho}$
  that is de Rham at $p$ with Hodge-Tate weights $0 \leq h_1 < h_2$ and $h_2 - h_1 \geq 1$
  is modular.
\end{theorem}

The key innovation of Kisin is the introduction of \emph{Kisin modules} (also called
$\mathfrak{S}$-modules or Breuil-Kisin modules):

\begin{definition}[Kisin module / Breuil-Kisin module]
  Let $K/\Qp$ be a finite totally ramified extension with uniformizer $\pi$ satisfying
  $E(\pi) = 0$ for an Eisenstein polynomial $E$. A \emph{Kisin module of height $\leq h$}
  is a finite free $\mathfrak{S} = W(k)[[u]]$-module $\mathfrak{M}$ (where $W(k)$ is the
  Witt ring of the residue field) together with a $\phi$-semilinear map
  $\phi_\mathfrak{M}\colon \mathfrak{M} \to \mathfrak{M}$ whose cokernel is killed by $E(u)^h$.
\end{definition}

\begin{theorem}[Kisin \cite{Kisin2006}]
  The functor $\mathfrak{M} \mapsto T(\mathfrak{M})$ from Kisin modules of height $\leq h$
  to $G_\infty$-stable $\Zp$-lattices (where $G_\infty = \Gal(\overline{K}/K_\infty)$ and
  $K_\infty = K(\pi^{1/p^\infty})$) is fully faithful and essentially surjective onto
  lattices in semistable representations of Hodge-Tate weights in $\{0,\ldots,h\}$.
\end{theorem}

% ==============================================================================
\section{Emerton's Completed Cohomology}\label{sec:emerton}
% ==============================================================================

Emerton \cite{Emerton2010} developed a new approach to the Fontaine-Mazur conjecture
using the $p$-adic local Langlands correspondence and \emph{completed cohomology}.

\begin{definition}[Completed cohomology, Emerton \cite{Emerton2006}]
  Let $Y(\Gamma)$ be the modular curve for a congruence subgroup $\Gamma \subset \GL_2(\Z)$.
  The \emph{completed cohomology} is
  \[
    \widetilde{H}^i := \varprojlim_n \varinjlim_\Gamma H^i(Y(\Gamma), \Z/p^n\Z),
  \]
  which carries a $\GL_2(\mathbb{A}_f^p) \times G_{\Qp}$-action.
\end{definition}

\begin{theorem}[Emerton \cite{Emerton2010}]\label{thm:emerton}
  Let $p \geq 5$ and let $\bar{\rho}\colon \GQ \to \GL_2(\F_p)$ be odd, absolutely
  irreducible, and modular. Assume that $\bar{\rho}|_{G_{\Qp}}$ is ``generic'' (a technical
  condition on the restriction to $G_{\Qp}$). Then any odd, irreducible, geometric
  $\rho\colon \GQ \to \GL_2(\Qp)$ with residual representation $\bar{\rho}$ is modular.
\end{theorem}

The strategy relies on Colmez's $p$-adic local Langlands functor and the observation that
if $\rho$ is geometric, then the representation $\Pi(\rho|_{G_{\Qp}})$ appears in
the completed cohomology, which implies the existence of a corresponding automorphic
representation.

\begin{remark}
  Together, Kisin's Theorem~\ref{thm:kisin} and Emerton's Theorem~\ref{thm:emerton} cover
  the vast majority of odd, irreducible, geometric $2$-dimensional $p$-adic Galois
  representations (subject to the genericity / non-exceptional conditions).
  The remaining cases (e.g., $p=2$, or representations with exceptional local behavior at
  $p$) are being addressed in ongoing work.
\end{remark}

% ==============================================================================
\section{Taylor's Potential Modularity}\label{sec:taylor}
% ==============================================================================

Even before full modularity was established, Taylor \cite{Taylor2002,Taylor2004}
proved a weaker but highly useful result: \emph{potential modularity}.

\begin{theorem}[Taylor \cite{Taylor2002}]\label{thm:taylor}
  Let $\ell$ be a prime and $\rho\colon \GQ \to \GL_2(\overline{\Q}_\ell)$ be an odd,
  irreducible, geometric $\ell$-adic representation with Hodge-Tate weights $\{0, k-1\}$,
  $k \geq 2$. Then there exists a totally real number field $F$ such that $\rho|_{G_F}$
  is modular (i.e., corresponds to a Hilbert modular form over $F$).
\end{theorem}

\begin{remark}
  Potential modularity suffices for many arithmetic applications, including:
  \begin{itemize}
    \item the meromorphic continuation and functional equation of $L(\rho, s)$;
    \item the Ramanujan conjecture for $\rho$ (Frobenius eigenvalues of absolute value
      $p^{(k-1)/2}$);
    \item cases of the Sato-Tate conjecture (now a theorem \cite{BGHT2011}).
  \end{itemize}
\end{remark}

\begin{theorem}[Barnet-Lamb-Gee-Geraghty-Taylor \cite{BGGT2014}]
  The Sato-Tate conjecture holds for all elliptic curves over totally real fields: the
  distribution of normalized Frobenius eigenvalues of $V_p E$ follows the Sato-Tate
  distribution (semicircle law on $[-2,2]$).
\end{theorem}

Taylor's method of ``tensor product tricks'' and the use of $\GL_n \times \GL_n$ for
suitable $n$ allows one to find a finite field extension over which the representation
becomes modular.

% ==============================================================================
\section{The $n \geq 3$ Case and Open Problems}\label{sec:openproblems}
% ==============================================================================

For $n \geq 3$, the Fontaine-Mazur conjecture remains almost entirely open.

\begin{conjecture}[Fontaine-Mazur for $n \geq 3$]\label{conj:FMn3}
  Every odd, irreducible, geometric $n$-dimensional $p$-adic Galois representation
  $\rho\colon \GQ \to \GL_n(\Qp)$ is automorphic: it corresponds to a cuspidal
  automorphic representation of $\GL_n(\mathbb{A}_\Q)$ via the Langlands correspondence.
\end{conjecture}

The primary obstacles are:
\begin{enumerate}[label=(\roman*)]
  \item \textbf{No $p$-adic local Langlands for $\GL_n$, $n \geq 3$}: Colmez's
    construction is intrinsically $\GL_2$-specific. The $p$-adic local Langlands for
    $\GL_3(\Qp)$ and beyond remains conjectural.
  \item \textbf{$R = \mathbb{T}$ for $\GL_n$}: The Taylor-Wiles method requires the
    ``numerical coincidence'' that deformation rings are isomorphic to Hecke rings, which
    requires a balance of local conditions that works cleanly for $\GL_2$ but faces
    obstructions for $\GL_n$.
  \item \textbf{Automorphic forms for $\GL_n$}: The theory of automorphic forms on
    $\GL_n$ is less concrete; in particular, there is no analogue of Shimura curves for
    $n \geq 3$ in the split case over $\Q$.
\end{enumerate}

\begin{remark}
  There is significant recent progress using \emph{perfectoid geometry} (Scholze
  \cite{Scholze2015}), \emph{prismatic cohomology} (Bhatt-Scholze \cite{BS2022}), and
  the categorical $p$-adic Langlands program (Fargues-Scholze \cite{FS2021}).
  These may eventually yield the tools to approach the $n \geq 3$ case.
\end{remark}

\begin{conjecture}[Categorical $p$-adic Langlands, Fargues-Scholze \cite{FS2021}]
  There is an equivalence of $\infty$-categories (after suitable localization) between
  \begin{itemize}
    \item coherent sheaves on the stack of Langlands parameters for $\GL_n(\Qp)$ over the
      Fargues-Fontaine curve, and
    \item smooth representations of $\GL_n(\Qp)$ in a suitable sense.
  \end{itemize}
\end{conjecture}

% ==============================================================================
\section{Known Results: Summary}
% ==============================================================================

\begin{theorem}[Status of Fontaine-Mazur for $n=2$]
  The following cases of Conjecture~\ref{conj:FM2} are proven:
  \begin{enumerate}[label=(\roman*)]
    \item All elliptic curves $E/\Q$: $V_pE$ is modular \cite{BCDT2001}.
    \item Odd, irreducible, $p$-adic representations $\rho$ with Hodge-Tate weights
      in $\{0,1\}$ and residually modular, under mild local conditions at $p$
      \cite{Kisin2009, Emerton2010}.
    \item All odd, irreducible, geometric $2$-dimensional representations become modular
      over some totally real field \cite{Taylor2002}.
    \item Weight $1$ forms: weight $1$ modular forms correspond to odd Artin
      representations, proven by Deligne-Serre \cite{DS1974}.
  \end{enumerate}
\end{theorem}

% ==============================================================================
\section{Conclusion}\label{sec:conclusion}
% ==============================================================================

The Fontaine-Mazur conjecture stands as one of the central open problems in arithmetic
geometry, unifying $p$-adic Hodge theory, Galois representations, and the Langlands
program. The $n=2$ case is largely settled through the combined work of Wiles, Kisin, and
Emerton, establishing that essentially all odd irreducible geometric $2$-dimensional
$p$-adic representations are modular. For $n \geq 3$, the conjecture remains open and
represents a major frontier of current research, with promising new tools emerging from
perfectoid geometry and the Fargues-Fontaine curve.

The resolution of the full Fontaine-Mazur conjecture would constitute a major step
towards the global Langlands correspondence and would imply, among other things, the
automorphic nature of all $\ell$-adic cohomology groups of smooth proper varieties over $\Q$.

\begin{thebibliography}{99}

\bibitem{BCDT2001}
C.\ Breuil, B.\ Conrad, F.\ Diamond, R.\ Taylor,
\textit{On the modularity of elliptic curves over $\mathbf{Q}$: wild $3$-adic exercises},
J.\ Amer.\ Math.\ Soc.\ \textbf{14} (2001), no.\ 4, 843--939.

\bibitem{Berger2002}
L.\ Berger,
\textit{Représentations $p$-adiques et équations différentielles},
Invent.\ Math.\ \textbf{148} (2002), no.\ 2, 219--284.

\bibitem{BGHT2011}
T.\ Barnet-Lamb, D.\ Geraghty, M.\ Harris, R.\ Taylor,
\textit{A family of Calabi-Yau varieties and potential automorphy II},
P.\ R.\ I.\ M.\ E.\ S.\ \textbf{47} (2011), no.\ 1, 29--98.

\bibitem{BGGT2014}
T.\ Barnet-Lamb, T.\ Gee, D.\ Geraghty, R.\ Taylor,
\textit{Potential automorphy and change of weight},
Ann.\ of Math.\ (2) \textbf{179} (2014), no.\ 2, 501--609.

\bibitem{BS2022}
B.\ Bhatt, P.\ Scholze,
\textit{Prisms and prismatic cohomology},
Ann.\ of Math.\ (2) \textbf{196} (2022), no.\ 3, 1135--1275.

\bibitem{Colmez2010}
P.\ Colmez,
\textit{Représentations de $\mathrm{GL}_2(\mathbf{Q}_p)$ et $(\phi, \Gamma)$-modules},
Astérisque \textbf{330} (2010), 281--509.

\bibitem{Deligne1971}
P.\ Deligne,
\textit{Formes modulaires et représentations $\ell$-adiques},
Séminaire Bourbaki, vol.\ 1968/69, Exposé 355, Lect.\ Notes in Math.\ vol.\ 179,
Springer (1971), 139--172.

\bibitem{DS1974}
P.\ Deligne, J.-P.\ Serre,
\textit{Formes modulaires de poids $1$},
Ann.\ Sci.\ École Norm.\ Sup.\ (4) \textbf{7} (1974), 507--530.

\bibitem{Emerton2006}
M.\ Emerton,
\textit{On the interpolation of systems of eigenvalues attached to automorphic Hecke eigenforms},
Invent.\ Math.\ \textbf{164} (2006), no.\ 1, 1--84.

\bibitem{Emerton2010}
M.\ Emerton,
\textit{Local-global compatibility in the $p$-adic Langlands programme for $\mathrm{GL}_{2/\mathbf{Q}}$},
preprint (2010), available at \url{https://math.uchicago.edu/~emerton/}.

\bibitem{FM1995}
J.-M.\ Fontaine, B.\ Mazur,
\textit{Geometric Galois representations},
in \textit{Elliptic Curves, Modular Forms, \& Fermat's Last Theorem (Hong Kong, 1993)},
Ser.\ Number Theory I, Int.\ Press, Cambridge, MA (1995), 41--78.

\bibitem{Fontaine1979}
J.-M.\ Fontaine,
\textit{Modules galoisiens, modules filtrés et anneaux de Barsotti-Tate},
in \textit{Journées de géométrie algébrique de Rennes}, Astérisque \textbf{65} (1979), 3--80.

\bibitem{Fontaine1982}
J.-M.\ Fontaine,
\textit{Sur certains types de représentations $p$-adiques du groupe de Galois d'un corps local;
construction d'un anneau de Barsotti-Tate},
Ann.\ of Math.\ (2) \textbf{115} (1982), no.\ 3, 529--577.

\bibitem{Fontaine1994}
J.-M.\ Fontaine,
\textit{Le corps des périodes $p$-adiques},
in \textit{Périodes $p$-adiques}, Astérisque \textbf{223} (1994), 59--111.

\bibitem{FS2021}
L.\ Fargues, P.\ Scholze,
\textit{Geometrization of the local Langlands correspondence},
preprint (2021), arXiv:2102.13459.

\bibitem{HT2001}
M.\ Harris, R.\ Taylor,
\textit{The Geometry and Cohomology of Some Simple Shimura Varieties},
Ann.\ of Math.\ Studies, vol.\ 151, Princeton Univ.\ Press (2001).

\bibitem{Henniart2000}
G.\ Henniart,
\textit{Une preuve simple des conjectures de Langlands pour $\mathrm{GL}(n)$ sur un corps $p$-adique},
Invent.\ Math.\ \textbf{139} (2000), no.\ 2, 439--455.

\bibitem{Kisin2006}
M.\ Kisin,
\textit{Crystalline representations and $F$-crystals},
in \textit{Algebraic Geometry and Number Theory}, Progr.\ Math.\ \textbf{253},
Birkhäuser (2006), 459--496.

\bibitem{Kisin2009}
M.\ Kisin,
\textit{Moduli of finite flat group schemes, and modularity},
Ann.\ of Math.\ (2) \textbf{170} (2009), no.\ 3, 1085--1180.

\bibitem{Scholze2015}
P.\ Scholze,
\textit{On torsion in the cohomology of locally symmetric varieties},
Ann.\ of Math.\ (2) \textbf{182} (2015), no.\ 3, 945--1066.

\bibitem{Taylor2002}
R.\ Taylor,
\textit{Remarks on a conjecture of Fontaine and Mazur},
J.\ Inst.\ Math.\ Jussieu \textbf{1} (2002), no.\ 1, 125--143.

\bibitem{Taylor2004}
R.\ Taylor,
\textit{Galois representations},
Ann.\ Fac.\ Sci.\ Toulouse Math.\ (6) \textbf{13} (2004), no.\ 1, 73--119.

\bibitem{TW1995}
R.\ Taylor, A.\ Wiles,
\textit{Ring-theoretic properties of certain Hecke algebras},
Ann.\ of Math.\ (2) \textbf{141} (1995), no.\ 3, 553--572.

\bibitem{Wiles1995}
A.\ Wiles,
\textit{Modular elliptic curves and Fermat's last theorem},
Ann.\ of Math.\ (2) \textbf{141} (1995), no.\ 3, 443--551.

\end{thebibliography}

\end{document}
